% theme: volcanoes_and_igneous_rocks
\MajorQuestion{火山と火成岩}
\SetRowSpaceQanda

\Instruction{地下にある物質と、噴火で出るものに関する次の問いに答えなさい。}
\QQ{地下で、岩石が熱などによってどろどろにとけたものを何というか。}{}
\QQ{地下にあったどろどろの物質が地表に流れ出たものを何というか。}{}
\QQ{火山が噴火したときに噴き出されたものをまとめて何というか。}{}
\QQ{噴火のときに噴き出す気体を何というか。}{}
\QQ{噴火で噴き出す、直径2\,mm以下の細かい粒を何というか。}{}

\Instruction{火山の形や噴火のようす、防災に関する次の問いに答えなさい。}
\QQ{傾斜がゆるやかで、平たく広がった形の火山を何というか。}{}
\QQ{溶岩と火山灰などが交互に積み重なった、円錐形の火山を何というか。}{}
\QQ{溶岩が火口付近に盛り上がってできた、ドーム状の火山を何というか。}{}
\QQ{噴火のようすや火山の形を左右する、どろどろした物質の流れにくさを何というか。}{}
\QQ{予想される災害の範囲や避難場所などを示した地図を何というか。}{}

\Instruction{地下にある物質が冷えてできる岩石に関する次の問いに答えなさい。}
\QQ{地下にあるどろどろした物質が冷えて固まった岩石を何というか。}{}
\QQ{地表や地表付近で急に冷えて固まった岩石を何というか。}{}
\QQ{地下深くでゆっくり冷えて固まった岩石を何というか。}{}
\QQ{大きな結晶と、そのまわりの細かい部分からなるつくりを何というか。}{}
\QQ{ほぼ同じ大きさの結晶が組み合わさったつくりを何というか。}{}

\Instruction{岩石に見られる粒と、代表的な岩石に関する次の問いに答えなさい。}
\QQ{地下にあるどろどろした物質の中でできる、結晶の粒を何というか。}{}
\QQ{急に冷えて固まった岩石に見られる、大きな結晶を何というか。}{}
\QQ{急に冷えて固まった岩石で、大きな結晶のまわりを埋める細かい部分を何というか。}{}
\QQ{黒っぽい色をした代表的な火山岩を何というか。}{}
\QQ{白っぽい色をした代表的な深成岩を何というか。}{}
